\chapter{Connected descent and seam corings}
\label{v4-ch:connected-descent}

Over a field \(k\), relative coefficient envelopes retain the coefficient algebra at the empty
configuration. A connected envelope instead has empty coefficient \(k\).
This change is already visible when two relative regular interfaces are
restricted to a connected envelope and composed:
their empty coefficient is \(B\otimes_k B\), and comparison with composition
over the relative envelope applies multiplication
\[
 B\otimes_k B\longrightarrow B.
\]
For \(B=\operatorname{Mat}_2(k)\), the source has dimension sixteen and the
target dimension four. The twelve-dimensional kernel survives although
every \(B\)-module is projective.

The object that records this difference is an extension--restriction coring.
It occurs at every junction of restricted interfaces. We first construct
it and compute its counit defect. A separate matrix calculation then shows
that the failure cannot be removed by choosing a different strong unital
lift to the same connected envelopes. Finally, under explicit hypotheses
on the arity filtration, we reconstruct connected modules from the whole
coring coaction. Reconstruction uses all its iterates; multiplication alone
does not retain the required data.

\section{The coring of an algebra map}
\label{v4-sec:seam-coring}

Let \((\mathcal C,\star,K)\) be a stable presentable monoidal
\(\infty\)-category whose product preserves colimits separately in
each variable. Fix a unital associative algebra map
\[
 f:G\longrightarrow R
\]
in \(\mathcal C\). Module categories and tensor products are derived.
We write \(X\star_G Y\) for the realization of the two-sided bar of a
right \(G\)-module \(X\) and a left \(G\)-module \(Y\). Thus
\[
 X\star_G Y
 =
 \left|[n]\longmapsto X\star G^{\star n}\star Y\right|.
\]
No interchange of two geometric factors is part of this definition.

Define the \(R\)--\(R\) bimodule
\[
 T_f=R\star_G R.
\]
Its outer actions multiply the two outside \(R\)-factors. Multiplication
gives a bimodule map
\[
 \epsilon:T_f\longrightarrow R.
\]
To define a coproduct, insert \(f\) into the middle of
\(R\star_G G\star_G R\):
\begin{equation}\label{v4-eq:seam-coproduct}
\Delta:T_f\longrightarrow
 R\star_G R\star_G R
 \simeq T_f\star_R T_f.
\end{equation}
This map is called the \emph{comultiplication} of the coring.
On an ordinary flat representative, this is the familiar formula
\[
 r\otimes s\longmapsto
 (r\otimes1)\otimes_R(1\otimes s).
\]
Equation~\eqref{v4-eq:seam-coproduct}, rather than an underived tensor
formula, specifies the map in general.

\begin{proposition}\label{v4-prop:seam-coring}
The maps \(\Delta\) and \(\epsilon\) make \(T_f\) a counital
coassociative coring in \(R\)-bimodules, with all higher coherences.
For the extension--restriction adjunction
\[
 F=R\star_G-:\operatorname{Mod}_G(\mathcal C)
 \ \rightleftarrows\
 \operatorname{Mod}_R(\mathcal C):U,
\]
the comonad \(FU\) is \(T_f\star_R-\).
\end{proposition}

\begin{proof}
Two successive insertions of \(f\) give the same ordered tensor
\(R\star_G R\star_G R\star_G R\). Applying \(\epsilon\) on either
side of \(\Delta\) removes one inserted unit and gives the identity.
For any number of insertions, the comparisons are induced by the same
multisimplicial bar. Its unit and associativity maps give the coherences
after realization.

For an \(R\)-module \(N\), associativity of the balanced tensor gives
\[
 FU(N)=R\star_G U(N)
 \simeq (R\star_G R)\star_R N.
\]
The adjunction counit is the action of \(R\), hence is induced by
\(\epsilon\). If \(\eta:\operatorname{Id}\to UF\) is the adjunction
unit, the comultiplication \(F\eta U\) inserts \(f\) and is induced by
\(\Delta\). These identify the comonad structure as well as its functor.
\end{proof}

A coring coaction on \(N\) is a map
\(N\to T_f\star_R N\) with its counit, coassociativity, and higher
compatibilities. We use this coherent notion throughout.
The coring need not have a bimodule coaugmentation
\(R\to T_f\). For example, at a scalar extension \(k\to B\), the
map \(b\mapsto b\otimes1\) is left \(B\)-linear but generally fails
right \(B\)-linearity.

\section{Composition after restriction}
\label{v4-sec:restricted-composition}

Let \(X\) be a right \(R\)-module and \(Y\) a left \(R\)-module.
We also write \(U\) for restriction of right modules.
Restriction gives a natural comparison
\[
 \chi_{X,Y}:U(X)\star_G U(Y)\longrightarrow X\star_R Y.
\]
It applies \(f\) to every intermediate \(G\)-factor in the two-sided bar.
Let \(D_f=\operatorname{cofib}(\epsilon:T_f\to R)\).

\begin{theorem}[Composition with its seam coring]
\label{v4-thm:seam-composition}
There are natural equivalences
\begin{align}
 U(X)\star_G U(Y)
 &\simeq X\star_R T_f\star_R Y,
 \label{v4-eq:seam-composition}\\
 \operatorname{cofib}(\chi_{X,Y})
 &\simeq X\star_R D_f\star_R Y.
 \label{v4-eq:seam-defect}
\end{align}
Under the first equivalence, \(\chi_{X,Y}\) applies the counit at the
middle factor. For a finite string of restricted bimodules, one copy
of the corresponding coring occurs at every intermediate algebra.
These identifications are compatible with regrouping the string.
\end{theorem}

\begin{proof}
Insert the regular \(R\)-modules next to \(X\) and \(Y\):
\[
 X\star_R(R\star_G R)\star_R Y
 \simeq
 (X\star_R R)\star_G(R\star_R Y)
 \simeq U(X)\star_G U(Y).
\]
The equivalences follow by realizing the common multibar, in which
all factors remain in their displayed order. Multiplication of the
middle two \(R\)-factors is exactly the bar comparison defining
\(\chi_{X,Y}\). Derived tensoring preserves cofibres, giving
\eqref{v4-eq:seam-defect}. The same multibar proves the assertion for
every finite string and all its parenthesizations.
\end{proof}

Apply this to the relative envelopes of
Chapter~\ref{v4-ch:relative-coefficient-envelopes}.
Suppose actual unital comparison maps
\[
 \theta_B:G_B\longrightarrow R_B=A\otimes_k B
\]
have been supplied for the coefficient algebras under consideration.
For a \(B\)--\(E\) bimodule \(P\), restrict
\(J_{B,E}(P)=A\otimes_k P\) along both endpoint maps and denote the
result by \(I^\theta_{B,E}(P)\).
For a composable \(E\)--\(D\) bimodule \(Q\),
\begin{equation}\label{v4-eq:coefficient-seam}
 I^\theta(P)\star_{G_E}I^\theta(Q)
 \simeq
 \operatorname{Res}\bigl(
 J(P)\star_{R_E}T_{\theta_E}\star_{R_E}J(Q)\bigr).
\end{equation}
The comparison to \(I^\theta(P\otimes_E^{\mathbb L}Q)\) applies the
counit and then Theorem~\ref{v4-thm:relative-composition}.
Its unit comparison is \(\theta_B:G_B\to R_B\).
Consequently restriction gives a lax composition law. For a selected
pair \(P,Q\), this law is an equivalence exactly when the right-hand
side of \eqref{v4-eq:seam-defect} vanishes for \(X=J(P)\), \(Y=J(Q)\).
A nonzero \(D_{\theta_E}\) does not by itself decide every such pair.

\section{The empty coefficient and noncommutative differentials}
\label{v4-sec:empty-seam}

For this section, let \(k\) be a field. Assume a colimit-preserving
strong monoidal functor \(e_0:\mathcal C\to\mathcal D(k)\) is given, and that
\[
 e_0(G)=k,\qquad e_0(R)=B,\qquad e_0(f)=(k\longrightarrow B),
\]
where \(B\) is an ordinary nonzero unital \(k\)-algebra.
Since \(e_0\) preserves the balanced bar,
\[
 e_0(T_f)=B\otimes_k B,\qquad e_0(\epsilon)=\mu_B.
\]
Here \(\mu_B(b\otimes c)=bc\).

Define the bimodule of noncommutative one-forms by
\[
 \Omega^1_{\mathrm{nc}}(B/k)
 =\ker\bigl(\mu_B:B\otimes_k B\to B\bigr),
\]
with outer \(B\)-actions, and put
\[
 d_{\mathrm{nc}}b=1\otimes b-b\otimes1.
\]
Commutativity of \(B\) does not require the two actions on this bimodule
to coincide; it does not identify it with the K\"ahler differential module.

\begin{proposition}\label{v4-prop:empty-seam-defect}
There is an equivalence
\[
 e_0(D_f)\simeq \Omega^1_{\mathrm{nc}}(B/k)[1].
\]
The kernel is zero exactly when the unit \(k\to B\) is an isomorphism.
If \(\dim_kB=r<\infty\), then
\[
 \dim_k\Omega^1_{\mathrm{nc}}(B/k)=r^2-r.
\]
In particular this dimension is \(d^4-d^2\) for
\(B=\operatorname{Mat}_d(k)\).
\end{proposition}

\begin{proof}
Multiplication has the \(k\)-linear section \(b\mapsto1\otimes b\).
Its cofibre is therefore its kernel shifted by one, and its rank is
\(\dim_kB\). This proves the equivalence and the finite-dimensional
formula.

Choose a \(k\)-linear functional \(\lambda:B\to k\) with
\(\lambda(1)=1\). If the kernel vanishes, then
\(1\otimes b-b\otimes1=0\) for every \(b\).
Applying \(\lambda\otimes\operatorname{id}\) gives \(b=\lambda(b)1\).
Thus \(B=k1\). The converse follows from multiplication in \(k\).
\end{proof}

The universal property explains which information this kernel carries.
For a derivation \(D:B\to M\) into a \(B\)-bimodule, define
\[
 h_D\left(\sum_i a_i\otimes b_i\right)=\sum_i a_iD(b_i)
 \quad\text{when }\sum_i a_ib_i=0.
\]
The derivation identity makes \(h_D\) right \(B\)-linear; left
linearity is immediate. Moreover
\(h_D(d_{\mathrm{nc}}b)=D(b)\), and
\[
 \sum_i a_i\,d_{\mathrm{nc}}b_i
 =\sum_i a_i\otimes b_i
\]
on the kernel. These formulas prove existence and uniqueness of the
factorization \(D=h_Dd_{\mathrm{nc}}\).
For a nonscalar \(b\), the tensor \(d_{\mathrm{nc}}b\) is nonzero:
a functional vanishing on \(1\) and taking \(b\) to \(1\), applied
to its first factor, gives \(-1\).

Thus \(B\ne k\) prevents the global counit \(T_f\to R\) from being
an equivalence. The converse requires information about the positive
configurations. In a family for which \(G_k=R_k\) and
\(\theta_k=\operatorname{id}\), it holds: the counit is then
\(R_k\star_{R_k}R_k\simeq R_k\). Empty evaluation alone proves no
such converse for an arbitrary map \(G\to R\).

\section{A matrix obstruction to strong connected transport}
\label{v4-sec:connected-Morita-obstruction}

The preceding defect concerns the specified restriction along \(f\).
We now test whether some other assignment could preserve all
coefficient bimodule compositions while keeping the connected
envelopes.

Let \(k\) have characteristic zero, let \(X=\mathbb A^1_k\), and
write \(\mathcal D(X)\) for the derived category of right differential
modules on \(X\), enhanced as a stable \(\infty\)-category.
Consider
\[
 \mathcal T_X=\mathcal D(k)\times\mathcal D(X)
\]
with symmetric monoidal product
\begin{equation}\label{v4-eq:square-zero-truncation}
 (U,M)\odot(V,N)
 =
 (U\otimes_k V,\ U\otimes_k N\oplus M\otimes_k V).
\end{equation}
The two components have arities zero and one. Two arity-one
factors have zero product in this category. The symmetry only
moves scalar complexes past differential modules, with Koszul signs.
It is not a statement that an intrinsic tensor product of two
differential modules vanishes.

Suppose a colimit-preserving strong monoidal truncation
\(\tau_{\leq1}:\mathcal C\to\mathcal T_X\) is given.
It sends balanced bars to balanced bars and therefore preserves
inverse bimodule interfaces.

For a unital algebra \(S\) in \(\mathcal T_X\), let
\(\operatorname{CH}(S;L)\) be its Hochschild chain object with
coefficients in an \(S\)-bimodule \(L\).
It is the realization of the cyclic bar with term
\(L\odot S^{\odot n}\) in degree \(n\). The last face uses the
symmetry to act at the left endpoint.

\begin{lemma}\label{v4-lem:singleton-chains}
For \(S=(k,M)\), with its scalar unit and actions, the arity-one
component is
\[
 \operatorname{CH}(S;S)_1\simeq M\oplus M[1].
\]
\end{lemma}

\begin{proof}
The unit splits, with reduced object \(\overline S=(0,M)\).
The normalized bar terms are
\(S\odot\overline S^{\odot n}[n]\).
At arity one only \(n=0,1\) survive. They are \(M\) and \(M[1]\),
and their boundary is \(1m-m1=0\).
The split unit-degenerate subcomplex is contractible. Thus the
normalization computes the full realization and retains the
internal differential on both copies of \(M\).
\end{proof}

\begin{lemma}\label{v4-lem:cyclic-two-bar}
For bimodules \({}_SP_T\) and \({}_TQ_S\) in \(\mathcal T_X\),
there is a natural equivalence
\[
 \operatorname{CH}(S;P\odot_T Q)
 \simeq
 \operatorname{CH}(T;Q\odot_S P).
\]
Consequently inverse bimodule interfaces preserve the Hochschild
chain object.
\end{lemma}

\begin{proof}
Resolve the balanced tensor and the Hochschild chains independently.
The resulting bisimplicial terms on the two sides are
\[
 P\odot T^{\odot q}\odot Q\odot S^{\odot p},
 \qquad
 Q\odot S^{\odot p}\odot P\odot T^{\odot q}.
\]
Rotate the two displayed blocks. On homogeneous tensors the sign is
\[
 (-1)^{(|p_0|+\sum_j|t_j|)(|q_0|+\sum_i|s_i|)},
\]
where \(p_0\in P\) and \(q_0\in Q\).
Interior faces multiply the same adjacent factors.
The two endpoint faces in the \(T\)-direction use the right action
on \(P\) and the left action on \(Q\); after rotation one becomes
the cyclic endpoint face. The \(S\)-direction uses the other two
actions in the same way. The Koszul symmetry supplies exactly
their endpoint signs, and inserted units have degree zero.
Thus the diagrams agree after their two indices are exchanged.

For the totalization convention
\eqref{v4-eq:bar-total}, exchanging degrees \(p,q\) adds
\((-1)^{pq}\). Indeed, their two bar differentials have signs
\(b_p+(-1)^p b_q\) before the exchange and
\(b_q+(-1)^q b_p\) afterwards.
The internal differential has sign \((-1)^{p+q}\) on both sides.
This verifies the chain comparison. Iterated realizations commute.
If \(P\odot_T Q\simeq S\) and \(Q\odot_S P\simeq T\),
substitution proves the last assertion.
\end{proof}

Let \(\omega_X\) denote the right differential module of top forms,
placed in cohomological degree zero.

\begin{theorem}\label{v4-thm:connected-Morita-obstruction}
Suppose connected envelopes \(G_B\), for ordinary unital
\(k\)-algebras \(B\), satisfy
\[
 \tau_{\leq1}(G_B)=(k,\omega_X\otimes_k B)
\]
with the indicated unit and scalar actions.
There is no assignment on all coefficient bimodules, their maps,
and equivalences that sends the regular bimodule \(B\) to \(G_B\)
and every derived coefficient composition to an equivalence of
composites over these envelopes.
The obstruction already occurs for \(\operatorname{Mat}_d(k)\)
and \(k\), with \(d\geq2\).
\end{theorem}

\begin{proof}
Set \(B=\operatorname{Mat}_d(k)\), \(e=E_{11}\),
\(P=Be\), and \(Q=eB\).
Matrix multiplication gives
\[
 P\otimes_k Q\simeq B,\qquad Q\otimes_B P\simeq k.
\]
For the first map the basis tensors
\(E_{i1}\otimes E_{1j}\) map to \(E_{ij}\).
For the second,
\(E_{1i}\otimes E_{j1}=\delta_{ij}(e\otimes e)\);
the inverse sends \(1\) to \(e\otimes e\).
All modules are projective on the balanced sides, so these also
compute the derived tensors. Associativity of matrix multiplication
gives their two Morita triangle identities.

The proposed assignment would yield inverse bimodule interfaces
between \(G_B\) and \(G_k\). Applying \(\tau_{\leq1}\) and
Lemma~\ref{v4-lem:cyclic-two-bar} would identify their Hochschild
chain objects. By Lemma~\ref{v4-lem:singleton-chains}, the arity-one
components would be
\[
 (\omega_X\otimes_k B)\oplus(\omega_X\otimes_k B)[1]
 \simeq \omega_X\oplus\omega_X[1].
\]
Their degree-zero cohomology differential modules have underlying
\(\mathcal O_X\)-ranks \(d^2\) and \(1\), which is impossible.
\end{proof}

This obstruction applies to every assignment with the stated
composition and unit requirements. It does not require its empty
component to equal the original bimodule.
It leaves the relative construction, lax restriction, and
restricted classes of coefficient interfaces available.

\section{Reconstruction along an arity filtration}
\label{v4-sec:arity-reconstruction}

Reconstruction asks whether the extended module
\(F(M)\), equipped with its canonical coring coaction, determines \(M\).
The answer below permits unbounded cohomological degrees. Its
convergence comes from arity.

For this section \(k\) is any field.
Assume that objects of \(\mathcal C\) have nonnegative arities and
that truncation to arities at most \(N\) gives exact strong monoidal
functors
\[
 q_N:\mathcal C\longrightarrow\mathcal C_{\leq N}.
\]
Assume \(\mathcal C_{\leq0}=\mathcal D(k)\), with \(K\) supported
at arity zero and \(K_0=k\).
Assume these functors preserve limits and colimits, jointly detect
equivalences, and have the same properties on the module and
bimodule categories in use. Products add lower bounds on arity:
a product of \(r\) positive-arity objects vanishes in arities
less than \(r\).
These are hypotheses on the enhanced carrier. They must be
verified before applying the theorem to a geometric coefficient
system.

Let \(G\to K\) be an algebra augmentation inducing \(G_0=k\) at
arity zero. Regard \(K\) as a \(G\)-bimodule through this
augmentation. Assume also that a \(G\)-bimodule map to \(K\)
is null whenever its arity-zero
component is null as a \(k\)-bimodule map.
This holds if extension by zero at arity zero is right adjoint
to arity-zero evaluation on \(G\)-bimodules.

Let \(f:G\to R\) be a unital algebra map whose arity-zero unit
\(k\to R_0\) splits in \(\mathcal D(k)\). Put
\begin{equation}\label{v4-eq:descent-fibres}
 I=\operatorname{fib}(G\to R),\qquad
 u:I\to G,\qquad
 J=\operatorname{fib}(G\to K).
\end{equation}
Here \(J\) has positive arity. In this section every power of
\(I\) or \(J\) is balanced over \(G\), in its written order.

\begin{lemma}\label{v4-lem:arity-nilpotence}
The map \(u\) factors through \(J\), and
\begin{equation}\label{v4-eq:arity-nilpotence}
 q_N\bigl(u^{\star_G(N+1)}:
 I^{\star_G(N+1)}\longrightarrow G\bigr)=0.
\end{equation}
\end{lemma}

\begin{proof}
The splitting of \(k\to R_0\) makes the fibre map \(I_0\to k\)
null. The assumption on maps to \(K\) therefore makes the composite
\(I\to G\to K\) null as a bimodule map. The fibre sequence
defining \(J\) gives the factorization.

Each term of the multibar for \(J^{\star_G r}\) contains \(r\)
positive-arity \(J\)-factors separated by nonnegative-arity
\(G\)-factors. It vanishes below arity \(r\).
Truncation preserves its realization, so
\(q_N(J^{\star_G(N+1)})=0\).
Tensoring the factorization of \(u\) with itself proves
\eqref{v4-eq:arity-nilpotence}.
\end{proof}

The nilpotence concerns a map. Its source need not vanish.
For example, for \(k\to\operatorname{Mat}_2(k)\) at arity zero,
\[
 I_0\simeq(\operatorname{Mat}_2(k)/k)[-1],\qquad
 \dim_k H^r(I_0^{\otimes r})=3^r .
\]
Every one of these tensor powers is nonzero, while the fibre map
to \(k\) is null.

For a left \(G\)-module \(M\), and \(r\geq1\), form the punctured
ordered cube
\begin{equation}\label{v4-eq:punctured-descent-cube}
 Q_r(M)=
 \lim_{\varnothing\ne S\subseteq\{1,\ldots,r\}}
 E_1(S)\star_G\cdots\star_G E_r(S)\star_G M,
 \qquad
 E_j(S)=
 \begin{cases}
 R,&j\in S,\\
 G,&j\notin S.
 \end{cases}
\end{equation}
The arrows insert \(f\). Restriction to subsets of
\(\{1,\ldots,r\}\) gives \(Q_{r+1}(M)\to Q_r(M)\).

\begin{theorem}[Ordered reconstruction]
\label{v4-thm:ordered-reconstruction}
Under the stated arity and splitting hypotheses, there are natural
fibre sequences
\[
 I^{\star_G r}\star_G M\longrightarrow M\longrightarrow Q_r(M),
\]
and a natural equivalence
\[
 M\simeq\lim_r Q_r(M).
\]
Consequently \(F=R\star_G-\) detects equivalences.
\end{theorem}

\begin{proof}
Adjoin the empty subset to the cube. Its initial vertex is \(M\).
Taking a fibre in each of its \(r\) directions replaces that
direction's \(G\to R\) by \(I\). Exactness of balanced tensor
therefore identifies the total fibre with
\(I^{\star_G r}\star_G M\). This proves the fibre sequence.

A transition on these fibres applies \(u\) to one terminal
\(I\)-factor. After \(q_N\), any \(N+1\) successive transitions
are null by \eqref{v4-eq:arity-nilpotence}.
The inverse limit of this fibre tower is zero.
Indeed, mapping any test object into it gives a tower of spectra
whose homotopy groups have uniformly zero transition composites.
Both the inverse limit and the first derived inverse limit of
each such group tower vanish. The Milnor exact sequence makes
the limit mapping spectrum contractible.
Preservation of limits by \(q_N\) and joint detection now give
the reconstruction equivalence.

If \(FM=0\), choose the rightmost \(R\)-factor in any nonempty
vertex of \eqref{v4-eq:punctured-descent-cube}.
All factors to its right are regular \(G\)-modules, so that
vertex contains \(R\star_G M=0\). Every \(Q_r(M)\) is zero,
and reconstruction gives \(M=0\).
Apply this to the fibre of a map to prove detection of
equivalences.
\end{proof}

There is no uniform nilpotence exponent in this argument:
the bound \(N+1\) is applied at fixed arity before taking the
arity limit. Nor does it bound coefficient resolution length.
That length can remain infinite even at arity zero.

\section{Effective coring descent}
\label{v4-sec:effective-coring-descent}

Conservativity of extension of scalars detects equivalences, but it does
not alone imply that every coherent coring coaction is effective.
The remaining condition is preservation of the totalizations
that become split after extension.

A split augmented cosimplicial object has an extra codegeneracy
contracting its augmentation. A tower is called \emph{pro-constant}
if its formal inverse system is equivalent to a constant object
in the pro-category. For towers, the mapping spaces there are
\[
 \operatorname{Map}_{\operatorname{Pro}}(\{X_i\},\{Y_j\})
 =
 \lim_j\operatorname*{colim}_i\operatorname{Map}(X_i,Y_j).
\]
Two facts will be used. Constant pro-objects in a stable
idempotent-complete category are closed under finite cofibres,
shifts, and retracts. Also, the partial-totalization tower of a
split augmented cosimplicial object is pro-constant at its
augmentation object. The latter follows because the split
augmentation remains a limit diagram under every functor;
applying the constant embedding into the pro-category gives
the assertion for the tower of finite totalizations.\footnote{See
A.~Mathew, \emph{The Galois group of a stable homotopy theory},
Example~3.11 and the proof of Proposition~3.22,
\href{https://arxiv.org/abs/1404.2156}{arXiv:1404.2156}.}

\begin{lemma}\label{v4-lem:split-totalizations}
At fixed arity at most \(N\), let \(Z^\bullet\) be a
cosimplicial \(G\)-module such that \(FZ^\bullet\) extends to a
split augmented cosimplicial \(R\)-module.
The partial-totalization tower of \(Z^\bullet\) is pro-constant.
Every exact functor preserves its totalization.
\end{lemma}

\begin{proof}
Work in the truncated module category and suppress \(q_N\).
Call a \(G\)--\(G\) bimodule \(P\) admissible if
\[
 \{\operatorname{Tot}_n(P\star_G Z^\bullet)\}_n
\]
is pro-constant. Since finite totalizations are finite stable
limits, admissible bimodules are closed under finite cofibres,
shifts, and retracts. The split augmentation makes \(R\)
admissible.

If \(P\) is admissible and \(Q\) is any bimodule,
the exact functor \(Q\star_G-\) preserves each finite
totalization. Its extension to pro-objects carries a constant
tower to a constant tower. Hence \(Q\star_G P\) is admissible.

Put \(d=N+1\). Factor the null map \(u^{\star_G d}\) through
\[
 I^{\star_G d}\longrightarrow I^{\star_G(d-1)}
 \longrightarrow\cdots\longrightarrow I\longrightarrow G,
\]
applying \(u\) on the right.
The successive cofibres are
\(I^{\star_G j}\star_G R\), \(0\leq j<d\).
They are admissible. The cofibre of the composite is a finite
extension of them, and \(G\) is a retract of this cofibre because
the composite is null. Thus \(G\) is admissible.

A pro-constant tower has inverse limit equal to its constant
value. An exact functor preserves finite totalizations and
carries that formal constant tower to the corresponding
constant tower of its image. It therefore preserves this
totalization, without needing to preserve arbitrary products.
\end{proof}

For an \(F\)-split cosimplicial object in the full category,
apply the lemma after each \(q_N\).
Since \(q_N\) preserves limits and balanced tensors, it
identifies the image of the comparison
\begin{equation}\label{v4-eq:descent-limit-comparison}
 F\operatorname{Tot}Z^\bullet
 \longrightarrow \operatorname{Tot}FZ^\bullet
\end{equation}
with the equivalence in the truncated category.
Joint detection proves \eqref{v4-eq:descent-limit-comparison}
is an equivalence in the full category.

Write \(K_f=FU=T_f\star_R-\). A coherent \(T_f\)-comodule means
a coalgebra for this comonad; denote its category by
\(\operatorname{Coalg}_{K_f}(\operatorname{Mod}_R(\mathcal C))\).

\begin{theorem}[Effective descent]
\label{v4-thm:effective-descent}
Under the hypotheses of Section~\ref{v4-sec:arity-reconstruction},
extension of scalars with its canonical coaction gives an equivalence
\[
 \operatorname{Mod}_G(\mathcal C)
 \simeq
 \operatorname{Coalg}_{K_f}
       \bigl(\operatorname{Mod}_R(\mathcal C)\bigr),
 \qquad
 M\longmapsto(FM,F\eta_M).
\]
For a coherent comodule \((N,\gamma)\), its inverse is
\begin{equation}\label{v4-eq:descent-inverse}
 \mathcal R_f(N,\gamma)
 =
 \operatorname{Tot}_{[n]\in\Delta}UK_f^nN.
\end{equation}
No cohomological boundedness assumption is imposed.
\end{theorem}

\begin{proof}
The adjunction supplies the comonad counit
\(\epsilon:K_f\to\operatorname{Id}\) and comultiplication
\(\delta=F\eta U\).
In the cosimplicial object of \eqref{v4-eq:descent-inverse},
the cofaces out of degree \(n\) are
\[
 d^0=\eta_{UK_f^nN},\qquad
 d^i=UK_f^{i-1}\delta_{K_f^{n-i}N}\ (1\leq i\leq n),
 \qquad d^{n+1}=UK_f^n\gamma.
\]
The codegeneracies apply the counit in each possible position.
Adjunction and coaction coherences give the entire cosimplicial
diagram.

Mapping a \(G\)-module into its totalization and applying the
adjunction at each vertex identifies the resulting mapping
space with maps of coherent \(K_f\)-coalgebras from
\((FM,F\eta_M)\) to \((N,\gamma)\).
Thus \(\mathcal R_f\) is right adjoint to the comparison functor.

After \(F\), the defining diagram has the split augmentation
\[
 N\xrightarrow{\gamma}K_fN
 \rightrightarrows K_f^2N\longrightarrow\cdots,
\]
whose extra codegeneracy is the counit at the left end.
The limit comparison already proved gives an equivalence
\[
 F\operatorname{Tot}UK_f^\bullet N
 \simeq \operatorname{Tot}K_f^{\bullet+1}N
 \simeq N .
\]
It respects coactions because it is induced by the full
cosimplicial diagram. This is the counit of the comparison
adjunction. It is therefore an equivalence.

The triangle identity shows that the comparison functor sends
its unit to an equivalence. The underlying functor of the
comparison is \(F\), which detects equivalences by
Theorem~\ref{v4-thm:ordered-reconstruction}.
Hence the unit is also an equivalence, proving the theorem.
\end{proof}

The totalization in \eqref{v4-eq:descent-inverse} is a limit.
Let \(\mathrm N^n\) denote cosimplicial normalization in degree \(n\).
In dg representatives the totalization uses the product
\[
 \prod_{n\geq0}\mathrm N^n(UK_f^\bullet N)[-n].
\]
This product belongs to the descent construction; it is not a
replacement of the direct-sum realizations used to define
balanced tensor products. The counit of the equivalence is the
map from this whole totalization. It is generally different
from the individual map \(K_fN\to N\).

Over a field, a nonacyclic unital dg algebra \(R_0\) has a split
unit as a complex. Indeed its unit cannot be a boundary: a
primitive for \(1\), multiplied by any cocycle, would make
every cocycle a boundary. Splitting complexes over the field
then supplies a chain map \(R_0\to k\) taking \(1\) to \(1\).
Thus, on a carrier satisfying the arity hypotheses and the
hypothesis on maps to \(K\), every such \(G\to R\) has effective
coring descent. Over a general arity-zero algebra \(S\), the
corresponding argument requires a splitting of \(S\to R_0\)
in \(\mathcal D(S\text{-}S)\). A splitting of underlying scalar complexes
does not supply that bimodule splitting.
